\scanpage{61}
et est $\Sigma$-exacte

(Variante: si $L\xrightarrow{u}M$ est \deleted{un $\Sigma$-épi} \inserted{\unsure{tel que|un}} \deleted{\unsure{alors il s'\illegible{1} aussi}} $\operatorname{Hom}(\Phi^i,u)$ est épi $\forall i$ (i.e. $u_{\bullet}:L_{\bullet}\longrightarrow M_{\bullet}$ épi) alors $u$ \deleted{\unsure{s'écrira}} admet un noyau (donc on a une suite $\Sigma$-exacte
\[
0\longrightarrow R\longrightarrow L\longrightarrow M\longrightarrow0\qquad)
\]

Utilisant ceci, on construit des objets de $P$ par ext. successives des objets $(\Psi_{\unsure{\varepsilon|i|\bullet}}^i)^{r_i}$ -- et on voit que $L\longmapsto L_{\bullet}=\operatorname{Hom}(\Phi^{\bullet},L)$ est un \uline{surjectif}.

e) Tout objet \inserted{$L$} de $P$ est extension successive d'objets de la forme $\Psi^i$ [qu'on peut supposer se suivre par ordre croissant, dans une filtration croissante \inserted{admissible} pour $P$ -- la quelle sera alors canonique.]

Il doit y avoir assez maintenant pour conclure que le foncteur $L\longmapsto L_{\bullet}=\operatorname{Hom}(\Phi^{\bullet},L)$ est ple. fid et a l'image essentielle les $L_{\bullet}$ tels que les $L_i$, $Z_i$, $B_i$, $H_i$ soient libres tf et \inserted{les $L_i$} nuls sauf un nb fini.
% archive sheet number: 60
